% =====================================================================
% experimento_solo_p_inter_costo_fixed.tex
% ---------------------------------------------------------------------
% Sección autocontenida: Experimento "solo_p_inter_20260531".
% Se integra en Capitulo_Experimentos.tex via \input{}.
% Sin preámbulo ni \begin{document}. Compilable con pdflatex >= 2 pases.
% =====================================================================

% Macro auxiliar de andamiaje: si ya existe en el documento padre se
% omite silenciosamente. Aquí se define de forma segura.
\providecommand{\TODO}[1]{\textcolor{rojomalo}{\textbf{[TODO: #1]}}}

% =====================================================================
\section{Experimento 8: solo\_p\_inter\_20260531 (primer ciclo con costo corregido)}
\label{sec:exp-solo-p-inter}
% =====================================================================

% =====================================================================
\subsection{Hipótesis y motivación}
% =====================================================================

El hallazgo documentado en la Sección~\ref{sec:evaluador-canonico}
reformuló el problema experimental. Los experimentos 1--6 operaron
sobre una función de costo que forzaba la orientación fija $u \to v$
para cada tarea, inflando sistemáticamente los gaps en $\approx 33$\,pp
y produciendo mínimos locales artificiales que resultaban idénticos
para las cinco metaheurísticas. La corrección consistió en integrar la
orientación \emph{greedy} (entrar a cada tarea por el extremo más
cercano al nodo previo) como la única lógica nativa de
\texttt{metacarp/evaluador\_costo.py}, con reporte de deadheading
consistente: \texttt{mejor\_costo} $\equiv$ \texttt{costo\_total\_desde\_reporte}.

Sobre esa base limpia, el nuevo programa experimental adopta el
principio metodológico \textbf{de lo más simple a lo más complejo}:
si el approach más básico ya produce gaps competitivos con el costo
corregido, los componentes avanzados (Path Relinking, kick reactivo,
AOS, presupuesto extendido) quedarán justificados únicamente como
reductores de tiempo de ejecución, no como mejoras de calidad
intrínseca. El presente experimento es el primero de ese nuevo ciclo.

La hipótesis central es la siguiente:

\begin{quote}
\textit{El selector probabilístico nativo
\texttt{seleccionar\_grupo\_operadores\_inter\_intra} ---sin
monkey-patches, sin Path Relinking, sin kick reactivo--- ya es
suficiente para obtener gaps bajos cuando el evaluador es correcto.
Calibrar únicamente $p_{\text{inter}}$ y el parámetro libre de cada
metaheurística debería rendir al menos tan bien como los approaches
más complejos del primer ciclo.}
\end{quote}

% =====================================================================
\subsection{Diseño del approach}
% =====================================================================

El approach \texttt{solo\_p\_inter} utiliza exclusivamente el
mecanismo de selección que ya posee el núcleo del proyecto. En cada
iteración, la función
\texttt{seleccionar\_grupo\_operadores\_inter\_intra} (módulo
\texttt{metacarp/metaheuristicas\_utils.py}) decide el grupo de
operadores según el estado de la solución actual:

\begin{itemize}
\item Si la solución \textbf{viola capacidad} ($\text{viol} > 10^{-12}$),
  la probabilidad de proponer el grupo INTER-ruta es
  $\alpha_{\text{inter}} = 0.80$ (reparación agresiva, fija, no
  calibrada). El propósito es redistribuir carga entre vehículos y
  recuperar la factibilidad.
\item Si la solución es \textbf{factible}, la probabilidad de proponer
  el grupo INTER-ruta es $p_{\text{inter}}$ (knob calibrado; exploración
  moderada para escapar de mínimos locales intra-ruta).
\end{itemize}

\textbf{Dentro del grupo elegido}, la función \texttt{generar\_vecino}
selecciona el operador de forma \textbf{uniforme}
(\texttt{pesos\_operadores=None}): no hay AOS ni pesos adaptativos.

El diseño NO instala ningún monkey-patch: el selector base es
exactamente el que comparten las cinco metaheurísticas (SA, TS simple,
RTS, ABC simple, Cuckoo). Esto elimina toda la capa de infraestructura
de monkey-patching del primer ciclo, con dos consecuencias
metodológicas importantes:

\begin{enumerate}
\item La comparación entre MHs es completamente limpia: todas corren
  el mismo mecanismo de selección sin interferencia externa.
\item El approach es reproducible con una única llamada a la función
  \texttt{*\_desde\_instancia} de cada MH con los kwargs apropiados;
  no requiere importar ni ordenar parches.
\end{enumerate}

Los demás parámetros de cada metaheurística se dejan
\textbf{instance-aware} (valor \texttt{None}): temperatura inicial y
mínima en SA, tenure bounds e iteraciones en TS/RTS, límite de
abandono en ABC, número de nidos y pasos de Lévy en Cuckoo. La
penalización de capacidad usa el $\lambda$ por defecto del módulo
(\texttt{lambda\_penal\_capacidad\_por\_defecto}, $\approx 10\times$
la mediana del costo de arco de la instancia, con mínimo absoluto
de 10), pasando \texttt{lambda\_capacidad=None}.

\medskip
\noindent\textbf{Conjunto de operadores.} Se activan los 9 operadores
del conjunto \texttt{OPERADORES\_POPULARES}:

\begin{itemize}
\item Intra-ruta (3): \texttt{relocate\_intra}, \texttt{swap\_intra},
  \texttt{2opt\_intra}.
\item Inter-ruta (6): \texttt{relocate\_inter}, \texttt{swap\_inter},
  \texttt{2opt\_star}, \texttt{cross\_exchange}, \texttt{or\_opt\_2},
  \texttt{or\_opt\_3}.
\end{itemize}

El uso del conjunto completo contrasta deliberadamente con los
experimentos 1--7, que empleaban el subconjunto reducido de 5
operadores (\texttt{OPERADORES\_STRICT\_5}). Dado que el objetivo del
experimento es medir la línea de base más simple posible, se prefiere
no sesgar la elección de operadores; el análisis de operadores del
primer ciclo (Tabla~\ref{tab:poda-ranking}) está disponible para
interpretar los resultados a posteriori.

\medskip
\noindent\textbf{Nota sobre \texttt{alpha\_inter}.} SA, TS simple y
RTS exponen el parámetro \texttt{alpha\_inter} en su firma y lo usan
directamente; se fija en $0.80$ en todos los casos. ABC simple y
Cuckoo no exponen ese kwarg: internamente calculan la probabilidad
en estado infactible como $\max(p_{\text{inter}},\,0.80)$, lo que
equivale a $0.80$ para todos los valores de $p_{\text{inter}}$
explorados en la malla ($\le 0.50$).

% =====================================================================
\subsection{Selector \texorpdfstring{$p_{\text{inter}}$}{p\_inter}: pseudocódigo}
% =====================================================================

El Algoritmo~\ref{alg:selector-solo-p-inter} formaliza el mecanismo.
Es idéntico al del Experimento 5 (Algoritmo~\ref{alg:selector-pip}),
pero aquí \textbf{no se instala como monkey-patch}: es la lógica
nativa del módulo, invocada por las cinco metaheurísticas en cada
iteración sin modificación.

\begin{algorithm}[h]
\caption{Selector $p_{\text{inter}}$ probabilístico nativo
         (\texttt{seleccionar\_grupo\_operadores\_inter\_intra}).}
\label{alg:selector-solo-p-inter}
\begin{algorithmic}[1]
\Require $\viol \in \mathbb{R}_{\ge 0}$ violación de capacidad actual;
  $p_{\text{inter}},\,\alpha_{\text{inter}} \in [0,1]$;
  generador $rng$;
  conjuntos $\mathcal{O}_{\text{intra}},\,\mathcal{O}_{\text{inter}},\,\mathcal{O}_{\text{fb}}$.
\Ensure Grupo de operadores $\mathcal{O}$ a usar en esta iteración.
\If{$\viol > 10^{-12}$}
  \State $p_{\text{ef}} \gets \alpha_{\text{inter}}$ \Comment{Estado infactible: reparación agresiva}
\Else
  \State $p_{\text{ef}} \gets p_{\text{inter}}$ \Comment{Estado factible: exploración moderada}
\EndIf
\If{$rng.\text{random}() < p_{\text{ef}}$ \textbf{y} $\mathcal{O}_{\text{inter}} \neq \emptyset$}
  \State \Return $\mathcal{O}_{\text{inter}}$
\ElsIf{$\mathcal{O}_{\text{intra}} \neq \emptyset$}
  \State \Return $\mathcal{O}_{\text{intra}}$
\Else
  \State \Return $\mathcal{O}_{\text{fb}}$ \Comment{Caso degenerado: grupos vacíos}
\EndIf
\end{algorithmic}
\end{algorithm}

Un dato relevante para la reproducibilidad: la función realiza
exactamente \textbf{un único} \texttt{rng.random()} por iteración.
Esto preserva la reproducibilidad bit a bit con la implementación
original del Recocido Simulado, de la que fue extraída.

% =====================================================================
\subsection{Grid search en dos fases}
% =====================================================================

El grid search sigue la misma lógica de las dos fases empleada en el
programa anterior, adaptada al nuevo approach.

\paragraph{Knob compartido.}
El único parámetro compartido entre las cinco metaheurísticas es
$p_{\text{inter}}$, que se calibra en la malla:
\[ p_{\text{inter}} \in \{0.1,\,0.2,\,0.3,\,0.4,\,0.5\}. \]

\paragraph{Parámetro libre por MH.}
Además del knob compartido, cada metaheurística expone un único
parámetro libre con tres valores posibles (el resto es instance-aware
o valor de literatura clásica):

\begin{table}[h]
\centering
\caption{Parámetros del grid por metaheurística. El resto de
         parámetros se dejan instance-aware (\texttt{None}).}
\label{tab:grid-solo-p-inter}
\small
\begin{tabular}{llll}
\toprule
MH & Parámetro libre & Valores & Resto notable \\
\midrule
SA           & \texttt{alpha} (enfriamiento)   & $\{0.90,\,0.95,\,0.99\}$  &
  $T_{\text{init}}, T_{\min}, L = n^2$: instance-aware \\
TS simple    & \texttt{tabu\_tenure}            & $\{7,\,15,\,25\}$         &
  tamaño vecindario: 25 \\
RTS          & \texttt{factor\_aumento}         & $\{1.1,\,1.2,\,1.3\}$    &
  factor\_reduccion: $0.9$ \\
ABC simple   & \texttt{num\_fuentes}            & $\{10,\,20,\,30\}$        &
  límite abandono: instance-aware \\
Cuckoo       & \texttt{pa\_abandono}            & $\{0.15,\,0.25,\,0.35\}$ &
  beta\_levy: $1.5$; nidos: instance-aware \\
\bottomrule
\end{tabular}
\end{table}

\noindent Esto produce $5 \times 3 = 15$ configuraciones por MH, para
un total de 75 configuraciones en el grid completo.

\paragraph{Fase 1 (calibración).}
Se ejecutan 3 repeticiones por (configuración, instancia) sobre las
23 instancias del corpus. No se fija semilla: la variabilidad es
puramente estocástica. Para cada MH se elige la configuración con el
\textbf{menor gap medio} sobre las 23 instancias. El resultado de esta
fase se escribe en \texttt{grid\_resumen.csv} (gap medio y desviación
estándar por configuración) y en \texttt{mejor\_config.json}
(configuración ganadora con su gap medio).

\paragraph{Fase 2 (corrida final).}
La configuración ganadora de la fase 1 se corre con \textbf{5
repeticiones} $\times$ 23 instancias. Los CSVs completos (con todas
las columnas de operadores y deadheading) se generan en el
subdirectorio \texttt{final/}.

\paragraph{Corridas totales.}
\begin{itemize}
\item Fase 1: $15\,\text{configs} \times 23\,\text{inst} \times
  3\,\text{reps} = 1{,}035$ corridas por MH
  $= 5{,}175$ corridas totales.
\item Fase 2: $23\,\text{inst} \times 5\,\text{reps} = 115$ corridas
  por MH $= 575$ corridas totales.
\item \textbf{Total}: 5{,}750 corridas entre ambas fases.
\end{itemize}

% =====================================================================
\subsection{Configuración y parámetros instance-aware}
% =====================================================================

La Tabla~\ref{tab:config-solo-p-inter} resume la configuración
completa del approach en comparación con los experimentos del primer
ciclo (evaluador canónico). Las diferencias clave se destacan con
negrita.

\begin{table}[h]
\centering
\caption{Comparativa de configuración entre el primer ciclo y
         \texttt{solo\_p\_inter\_20260531}.}
\label{tab:config-solo-p-inter}
\small
\begin{tabular}{lll}
\toprule
Componente & Primer ciclo (R1--R7) & \texttt{solo\_p\_inter} \\
\midrule
Evaluador de costo    & Canónico (artefacto)       & \textbf{Greedy nativo} \\
Selector inter/intra  & Monkey-patch (strict o PR) & \textbf{Selector nativo} \\
Path Relinking        & Sí (R3--R7)                & \textbf{No} \\
Kick reactivo         & Sí (R1--R7)                & \textbf{No} \\
AOS-PM                & Sí (R2)                    & \textbf{No} \\
Presupuesto extra     & $\times 25$ en R5--R7      & \textbf{Instance-aware} \\
$\lambda$             & Default / $\times 10$      & \textbf{Default (None)} \\
Conjunto operadores   & 5 (\texttt{STRICT\_5})     & \textbf{9 (\texttt{POPULARES})} \\
$\alpha_{\text{inter}}$ & $0.80$ (R4--R7)          & $0.80$ (fijo) \\
$p_{\text{inter}}$    & $0.20$ fijo (R4--R7)       & \textbf{calibrado en {[}0.1, 0.5{]}} \\
Reheat SA             & Default (sin tope)         & \textbf{Acotado} (\texttt{patience=10}, \texttt{max\_reheats=5}) \\
\bottomrule
\end{tabular}
\end{table}

\medskip
\noindent\textbf{Parámetros instance-aware.} Para cada metaheurística
los parámetros principales quedan determinados automáticamente por el
wrapper \texttt{*\_desde\_instancia} a partir del tamaño de la
instancia (número de tareas $n_\tau$, número de aristas, capacidad $Q$
y estadísticas del grafo):

\begin{itemize}
\item \textbf{SA}: $T_{\text{init}}$ derivada de la varianza del costo
  de vecinos muestreados; $T_{\min}$ como fracción de $T_{\text{init}}$;
  longitud de cadena $L = n_\tau^2$. Reheat acotado (véase nota
  abajo).
\item \textbf{TS simple / RTS}: número máximo de iteraciones y
  ventanas de tenure bounds derivados de $n_\tau$.
\item \textbf{ABC simple}: número de nidos derivado de $n_\tau$ (el
  \texttt{limite\_abandono} ya no es instance-aware: se calibró a $60$).
\item \textbf{Cuckoo}: número de nidos y pasos de vuelo de Lévy
  derivados de $n_\tau$.
\end{itemize}

\medskip
\noindent\textbf{Segundo parámetro calibrado por MH.} Además de
$p_{\text{inter}}$ y del knob principal, se calibró (sobre esta misma base
\texttt{solo\_p\_inter}, 3 reps $\times$ 23 instancias) el parámetro adicional
más influyente de cada MH, dejándolo fijo en su mejor valor: SA
\texttt{max\_reheats\_sin\_mejora}$=10$; TS \texttt{tam\_vecindario}$=40$; RTS
\texttt{factor\_reduccion}$=0.95$; ABC \texttt{limite\_abandono}$=60$; Cuckoo
\texttt{beta\_levy}$=1.3$. Las mayores ganancias fueron en TS y ABC.

% ---------------------------------------------------------------------
\subsubsection*{Nota técnica: reheat acotado en SA}
\label{nota:reheat-acotado-solo-p-inter}
% ---------------------------------------------------------------------

El Recocido Simulado utiliza un mecanismo de \emph{reheat} para
escapar de mínimos locales: cuando la temperatura cae demasiado sin
mejora del mejor global, se recalienta parcialmente. El comportamiento
por defecto del módulo SA fija \texttt{max\_reheats\_sin\_mejora=0}
(sin tope de reheats), lo que hace que corridas con \texttt{alpha}
alto ($\alpha = 0.99$, enfriamiento muy lento) no converjan en tiempo
razonable.

El approach \texttt{solo\_p\_inter} fija el reheat \textbf{acotado},
consistente con la configuración validada del proyecto:

\begin{itemize}
  \item \texttt{patience=10}: iteraciones de nivel sin mejora antes
    de activar el reheat.
  \item \texttt{reheat\_factor=0.5}: la temperatura de reheat es la
    mitad de la temperatura inicial efectiva.
  \item \texttt{max\_reheats\_sin\_mejora=10}: máximo de reheats
    consecutivos sin mejorar el mejor global antes de parar. Es el
    \emph{segundo knob} calibrado de SA (mejor de $\{3,5,10\}$).
\end{itemize}

Con este acotamiento, una corrida SA con $\alpha = 0.99$ tarda
$\sim$5\,s en instancias pequeñas, el mismo orden de magnitud que
$\alpha = 0.90$ ó $0.95$. Sin él, la misma corrida no termina en
150\,s (verificado experimentalmente). \texttt{patience} y
\texttt{reheat\_factor} se mantienen fijos; solo
\texttt{max\_reheats\_sin\_mejora} se calibró como segundo knob de SA.

El mismo reheat acotado se aplica en el Experimento~\ref{sec:exp-binario-capacidad}
(\texttt{binario\_capacidad\_20260531}), garantizando comparabilidad
directa entre ambos approaches.

% =====================================================================
\subsection{Reproducibilidad: scripts e infraestructura}
% =====================================================================

El experimento se organiza en tres niveles de scripts:

\begin{itemize}
\item \textbf{Orquestador común}:
  \texttt{scripts/\_solo\_p\_inter\_20260531\_common.py} --- implementa
  el grid search en dos fases, la construcción de kwargs por MH, la
  consolidación de CSV parciales y la CLI (\texttt{--smoke},
  \texttt{--solo-fase}, \texttt{--workers}).
\item \textbf{Runners por MH}:
  \texttt{scripts/run\_sa\_solo\_p\_inter\_20260531.py},
  \texttt{run\_tabu\_simple\_solo\_p\_inter\_20260531.py},
  \texttt{run\_tabu\_reactiva\_solo\_p\_inter\_20260531.py},
  \texttt{run\_abc\_simple\_solo\_p\_inter\_20260531.py},
  \texttt{run\_cuckoo\_solo\_p\_inter\_20260531.py} ---
  cada uno fija su MH y llama al orquestador.
\item \textbf{Lanzador master}:
  \texttt{scripts/run\_all\_solo\_p\_inter\_20260531.sh} --- lanza
  los cinco runners en paralelo (uno por proceso), con redirección de
  logs.
\end{itemize}

La salida se escribe en:
\begin{center}
\texttt{experimentos\_costo\_fixed/<mh>\_solo\_p\_inter\_<YYYYMMDD-HHMM>/}
\end{center}

Con la siguiente estructura de carpetas:

\begin{itemize}
\item \texttt{grid\_parciales/} --- CSVs parciales de la fase 1, un
  archivo por corrida (columnas completas con contadores de operadores
  y reporte de deadheading).
\item \texttt{calibracion\_todas.csv} --- consolidado de la fase 1
  (unión de columnas de todos los parciales).
\item \texttt{grid\_resumen.csv} --- gap medio y desviación estándar
  por configuración, ordenados de mejor a peor.
\item \texttt{mejor\_config.json} --- configuración ganadora con sus
  campos: \texttt{mh}, \texttt{p\_inter}, \texttt{libre\_nombre},
  \texttt{libre\_valor}, \texttt{gap\_medio}, \texttt{gap\_std}.
\item \texttt{final/} --- CSVs de la fase 2 (5 repeticiones, columnas
  completas), uno por instancia con el patrón
  \texttt{<mh>\_solo\_p\_inter\_<instancia>.csv}.
\end{itemize}

\medskip
\noindent\textbf{Columnas del CSV.} Las columnas son idénticas a las
del núcleo del proyecto. Entre las más relevantes para el análisis:
\texttt{gap\_bks\_porcentaje}, \texttt{mejor\_costo},
\texttt{costo\_total\_desde\_reporte},
\texttt{reporte\_detalle\_deadheading}, y los 36 contadores de
operadores ($9\text{ operadores} \times 4\text{ categorías}$:
propuesto, aceptado, mejoró, trayectoria\_mejor).

% =====================================================================
\subsection{Resultados}
% =====================================================================

\begin{observacion}[Corrida completada]
El experimento \texttt{solo\_p\_inter\_20260531} se ejecutó el
31 de mayo de 2026 sobre las 23 instancias, con grid search en dos
fases (3 repeticiones en la fase~1 de calibración y 5~repeticiones en
la fase~2 de confirmación). Los valores reportados a continuación
corresponden a la fase~2 con la configuración ganadora de cada MH.
\end{observacion}

\paragraph{Fuentes de datos.}
Cada número de la Tabla~\ref{tab:solo-p-inter-resultados} se extrae
de los siguientes archivos generados por el orquestador:

\begin{itemize}
\item \textbf{$p_{\text{inter}}^*$ y parámetro libre*}: columnas
  \texttt{p\_inter} y \texttt{<libre\_nombre>} de la primera fila
  de \texttt{grid\_resumen.csv} (o equivalentemente del campo
  correspondiente en \texttt{mejor\_config.json}).
\item \textbf{Gap medio (fase 2)}: promedio de la columna
  \texttt{gap\_bks\_porcentaje} en los CSV de
  \texttt{final/<mh>\_solo\_p\_inter\_<instancia>.csv} sobre las
  23 instancias $\times$ 5 repeticiones.
\item \textbf{$N$ instancias en BKS}: número de instancias para las
  que \texttt{gap\_bks\_porcentaje} $< 0.01$\,\% en al menos una de
  las 5 repeticiones de la fase 2.
\end{itemize}

\begin{table}[h]
\centering
\caption{Resultados del experimento \texttt{solo\_p\_inter\_20260531}
         (fase 2: 5 reps $\times$ 23 instancias por MH). El gap se mide
         respecto al BKS de cada instancia con el evaluador greedy nativo.}
\label{tab:solo-p-inter-resultados}
\small
\begin{tabular}{lcccccc}
\toprule
MH & $p_{\text{inter}}^*$ & Parámetro libre* & Valor* &
  Gap medio (\%) & Mediana (\%) & $N$ inst.\ en BKS \\
\midrule
SA           & 0.5 & \texttt{alpha}          & 0.90 &
  2.89  & 1.77  & 14/23 \\
TS simple    & 0.4 & \texttt{tabu\_tenure}   & 25   &
  5.76  & 3.64  & 8/23 \\
RTS          & 0.5 & \texttt{factor\_aumento}& 1.2  &
  4.97  & 3.64  & 10/23 \\
ABC simple   & 0.5 & \texttt{num\_fuentes}   & 30   &
  9.64  & 9.00  & 5/23 \\
Cuckoo       & 0.1 & \texttt{pa\_abandono}   & 0.15 &
  19.13 & 18.29 & 0/23 \\
\midrule
\textbf{Media} & --- & --- & --- & \textbf{8.48} & --- & 37/115 \\
\bottomrule
\end{tabular}
\\[4pt]
\footnotesize
* Configuración canónica fija (p\_inter, knob principal y segundo knob
calibrados; leída de \texttt{config\_fija.json}). El segundo knob calibrado por
MH es: SA \texttt{max\_reheats\_sin\_mejora}$=10$; TS
\texttt{tam\_vecindario}$=40$; RTS \texttt{factor\_reduccion}$=0.95$; ABC
\texttt{limite\_abandono}$=60$; Cuckoo \texttt{beta\_levy}$=1.3$. La fila
\textbf{Media} agrega los aciertos de BKS como pares (MH, instancia) sobre el
total $5 \times 23$ por MH.
\end{table}

\paragraph{Comparativa con R3 y R5 (referencia).}
La Tabla~\ref{tab:solo-p-inter-comparativa} sitúa los resultados de
este experimento respecto a los mejores approaches del primer ciclo
bajo el evaluador greedy (R3 Path Relinking y R5 Budget), extraídos
de la Tabla~\ref{tab:re-greedy-gap}.

\begin{table}[h]
\centering
\caption{Gap medio (\%) por MH: \texttt{solo\_p\_inter} (evaluador greedy nativo)
         frente a R3 y R5 (evaluador greedy por monkey-patch).
         $\Delta$ positivo indica que \texttt{solo\_p\_inter} es peor que R3.}
\label{tab:solo-p-inter-comparativa}
\small
\begin{tabular}{lrrrr}
\toprule
MH & R3 PR & R5 Budget & \texttt{solo\_p\_inter} & $\Delta$ vs R3 \\
\midrule
SA        & 0.90  & 0.57 & 2.89  & $+1.99$  \\
TS simple & 4.48  & 1.08 & 5.76  & $+1.28$  \\
RTS       & 3.04  & 1.18 & 4.97  & $+1.93$  \\
ABC       & 9.75  & 8.30 & 9.64  & $-0.11$  \\
Cuckoo    & 8.20  & 4.25 & 19.13 & $+10.93$ \\
\midrule
\textbf{Media} & 5.27 & 3.08 & \textbf{8.48} & $+3.21$ \\
\bottomrule
\end{tabular}
\end{table}

\paragraph{Interpretación de los resultados.}
Con \texttt{p\_inter}, el knob principal y el \emph{segundo} parámetro más
influyente de cada MH ya calibrados, el selector desnudo se acerca mucho a R3
(Path Relinking del primer ciclo): $\Delta$ medio $+3.21$\,pp y, en ABC, incluso
lo \textbf{supera} ($9.64$ vs $9.75$\,\%). Tres lecturas se desprenden:

\begin{enumerate}
\item \textbf{La calibración del segundo knob importa, sobre todo en TS y ABC.}
  Calibrar \texttt{tam\_vecindario}$=40$ (TS) y \texttt{limite\_abandono}$=60$
  (ABC) bajó sus gaps de $8.40$ a $5.76$\,\% y de $14.66$ a $9.64$\,\%
  respectivamente; en SA, RTS y Cuckoo el efecto fue marginal (sus segundos
  knobs ya estaban cerca del óptimo). La media del approach pasó de
  $10.21$ a $8.48$\,\%.
\item \textbf{El óptimo de $p_{\text{inter}}$ es alto.} La calibración
  selecciona $p_{\text{inter}} = 0.5$ en SA, RTS y ABC (y $0.4$ en TS),
  muy por encima del $0.20$ fijado en R4. Con el costo correcto, una
  mayor proporción de movimientos inter-ruta resulta beneficiosa, lo
  que sugiere que el valor moderado heredado del primer ciclo estaba
  sesgado por el artefacto de orientación. Cuckoo es la excepción
  ($p_{\text{inter}} = 0.1$).
\item \textbf{Jerarquía de MH como línea de base.} Con el approach más
  simple, SA ($2.89$\,\%) y RTS ($4.97$\,\%) son superiores a las
  metaheurísticas poblacionales (ABC $9.64$\,\%, Cuckoo $19.13$\,\%); ABC,
  tras calibrar su límite de abandono, se vuelve competitiva, mientras que
  Cuckoo sigue rezagada.
\end{enumerate}

En consecuencia, el resultado fija una \textbf{línea de base limpia} y bien
calibrada: todo mecanismo añadido en los siguientes approaches debe justificarse
por la mejora de gap que produzca \emph{sobre} estos valores, ya
medidos con el evaluador nativo. El siguiente experimento de control
(Experimento~\ref{sec:exp-binario-capacidad}) sustituye el selector
probabilístico por uno binario determinista para aislar el efecto del
mecanismo de selección antes de reintroducir la intensificación.

% =====================================================================
\subsection{Conclusión provisional del experimento 8}
% =====================================================================

Este experimento cierra el primer ciclo experimental y abre el segundo,
fundamentado en el evaluador greedy integrado de forma nativa. Más allá de los
gaps obtenidos, sus contribuciones metodológicas son:

\begin{enumerate}
\item \textbf{Línea de base limpia bajo el evaluador corregido.}
  Los resultados de \texttt{solo\_p\_inter} constituyen el primer
  punto de referencia genuino del programa experimental: sin artefactos
  de evaluación y sin componentes de complejidad artificial.
\item \textbf{Calibración por MH de $p_{\text{inter}}$ y de los dos knobs
  más influyentes.} La calibración cuantifica, bajo el evaluador correcto,
  el valor óptimo de $p_{\text{inter}}$, del knob principal y del segundo
  parámetro de cada metaheurística (p.ej. \texttt{tam\_vecindario} en TS,
  \texttt{limite\_abandono} en ABC). El resultado se fija como configuración
  canónica reutilizada por los demás approaches. Es independiente y
  complementario a los experimentos R4/R5 del primer ciclo (donde el evaluador
  greedy se aplicó como monkey-patch sobre una búsqueda ya diseñada con el
  evaluador canónico).
\item \textbf{Separación de responsabilidades.}
  Al no instalar monkey-patches, el experimento permite atribuir
  directamente la calidad del resultado a la metaheurística base y al
  evaluador, sin interferencia de la infraestructura de patching.
\end{enumerate}

Los resultados completos se incorporarán a la
Tabla~\ref{tab:solo-p-inter-resultados} una vez finalizada la corrida
en \texttt{experimentos\_costo\_fixed/}.
